% !TEX root = ../main.tex

\section{Silniki cieplne w termodynamice fenomenologicznej}
Silniki cieplne są maszynami konwertującymi cyklicznie dostarczane do nich ciepło na pracę mechaniczną, a czynnikiem roboczym jest płyn. Przebieg pełnego cyklu silnika cieplnego jest krzywą zamknięta na płaszczyźnie \(p\pab{V}\). Krzywa ta jest skierowana, co prowadzi do podzielenia cykli na prawobieżne i lewobieżne. Na rysunku \ref{fig:cykl_prawo_lewo} przedstawiono różnice pomiędzy cyklem prawobieżnym i lewobieżnym dla takiej samej krzywej zamkniętej. Podział ten odzwierciedla zasadnicze różnice w działaniu tych maszyn. Obiegi prawobieżne, których przykład przedstawiono na rysunku \ref{fig:cykl_prawobiezny}, wykonują pracę nad otoczeniem i są to silniki cieplne; obiegi lewobieżne, którego przykład przedstawiono na rysunku \ref{fig:cykl_lewobiezny}, powodują przepływ ciepła między zbiornikami kosztem wykonanej pracy nad układem i są to pompy ciepła. Pole ograniczone krzywą przebiegu reprezentuje wartość bezwzględną pracy objętościowej wykonanej w jednym cyklu.

Teoretyczny opis cykli termodynamicznych w przypadku ogólnym jest skomplikowany, ponieważ krzywa opisująca go może być dowolna. Aby je opisać wyróżnia się cykle wyidealizowane mające być przybliżeniem rzeczywistych przemian w silniku. Składają się one z elementarnych przemian termodynamicznych, które opisane są prostymi relacjami matematycznymi. Najważniejszym z cykli jest silnik Carnota, który ma duże znaczenie w termodynamice teoretycznej, a czynnikiem roboczym jest gaz doskonały.

\begin{figure}[H]
	\centering
	\begin{subfigure}{0.45\textwidth}
		\centering
		\input{figures/pv_silnik_prawo.tex}
		\caption{}
		\label{fig:cykl_prawobiezny}
	\end{subfigure}
	\begin{subfigure}{0.45\textwidth}
		\centering
		\tikzsetnextfilename{pv_silnik_lewo}
\begin{tikzpicture}
	\begin{axis}[width=\textwidth,
		xlabel=$V$,
		ylabel=$p$,
		axis x line=bottom,
		axis y line=middle,
		xmin=0,
		xmax=3,
		ymin=0,
		ymax=3,
		ticks=none,
		every axis x label/.style={at={(current axis.right of origin)},anchor=north east},
		every axis y label/.style={at={(current axis.above origin)},anchor=north east},
		declare function={
			centerx=1.5;
			centery=1.5;
			angle0=30;
			r = 1;
		}]
		\addplot[black,thick,->] ({centerx-r*cos(angle0)},{centery-r*sin(angle0)}) arc (180+angle0:270+angle0:r);
		\addplot[black,thick,->] ({centerx-r*sin(angle0)},{centery+r*cos(angle0)}) arc (90+angle0:180+angle0:r);
		\addplot[black,thick,->] ({centerx+r*cos(angle0)},{centery+r*sin(angle0)}) arc (angle0:90+angle0:r);
		\addplot[black,thick,->] ({centerx+r*sin(angle0)},{centery-r*cos(angle0)}) arc (-90+angle0:angle0:r);
	\end{axis}
\end{tikzpicture}

		\caption{}
		\label{fig:cykl_lewobiezny}
	\end{subfigure}
	\captionsetup{subrefformat=parens}
	\caption{Porównanie cyklu prawobieżnego \subref{fig:cykl_prawobiezny} i lewobieżnego \subref{fig:cykl_lewobiezny}.}
	\label{fig:cykl_prawo_lewo}
\end{figure}

W termodynamice technicznej powszechna jest konwencja, według której ciepło dostarczone z ciepłego zbiornika, odprowadzone do zimnego zbiornika oraz praca wykonana przez silnik są dodatnie \cite{b:huang_podstawy_fizyki_statystycznej}. Przepływy te zostały przedstawione na rysunku \ref{fig:konwencja_znakow_1}. Taki kierunek jest typowy podczas normalnej pracy. Tu mimo to użyta będzie konwencja powszechna w termodynamice statystycznej w celu ujednolicenia oznaczeń. W konwencji tej ciepła dostarczone do układu oraz praca wykonana nad układem mają znak dodatni. W przeciwnym przypadku mają one znak ujemny. Na rysunku \ref{fig:konwencja_znakow_2} przestawiono kierunku przepływu dodatniego.

\begin{figure}
	\centering
	\begin{subfigure}{0.45\textwidth}
		\centering
		\begin{tikzpicture}[x=0.5cm, y=0.5cm]
	\tikzmath{
		\tankwidth = 6;
		\tanksep = 5;
		\tankthickness = 1;
		\engwidth = 1;
		\engsep = 1;
		\radius = 2;
		\chansep = 1;
	}
	\draw[color=red, thick] (-\tankwidth, \tanksep) -- (\tankwidth, \tanksep);
	\draw[color=red, thick] (-\tankwidth, \tanksep) -- (-\tankwidth, \tanksep+\tankthickness);
	\draw[color=red, thick] (\tankwidth, \tanksep) -- (\tankwidth, \tanksep+\tankthickness);
	\fill[color=red, opacity=0.2] (-\tankwidth, \tanksep) rectangle (\tankwidth, \tanksep+\tankthickness);
	
	\draw[color=blue, thick] (-\tankwidth, -\tanksep) -- (\tankwidth, -\tanksep);
	\draw[color=blue, thick] (-\tankwidth, -\tanksep) -- (-\tankwidth, -\tanksep-\tankthickness);
	\draw[color=blue, thick] (\tankwidth, -\tanksep) -- (\tankwidth, -\tanksep-\tankthickness);
	\fill[color=blue, opacity=0.2] (-\tankwidth, -\tanksep) rectangle (\tankwidth, -\tanksep-\tankthickness);
	
	\draw[black] (-\engsep-\engwidth, -\tanksep) -- (-\engsep-\engwidth, \tanksep);
	\draw[black] (-\engsep+\engwidth, -\tanksep) -- (-\engsep+\engwidth, -\chansep-\radius);
	\draw[black] (-\engsep+\engwidth, \chansep+\radius) -- (-\engsep+\engwidth, \tanksep);
	\draw[black] (-\engsep+\engwidth, -\chansep-\radius) arc (180:90:\radius);
	\draw[black] (-\engsep+\engwidth, \chansep+\radius) arc (180:270:\radius);
	\node[anchor=east] at (-\engsep-\engwidth, 0.5*\tanksep)  {\(Q_h>0\)};
	\node[anchor=east] at (-\engsep-\engwidth, -0.5*\tanksep) {\(Q_c>0\)};
	\node[anchor=west] at (-\engsep+\engwidth+\radius, 0) {\(W>0\)};
	\draw[line width=0.3mm, -{Stealth[length=5mm, open]}] (-\engsep, 0.75*\tanksep) -- (-\engsep, 0.25*\tanksep);
	\draw[line width=0.3mm, -{Stealth[length=5mm, open]}] (-\engsep, -0.25*\tanksep) -- (-\engsep, -0.75*\tanksep);
	\draw[line width=0.3mm, -{Stealth[length=5mm, open]}] (-\engsep+\engwidth, 0) -- (-\engsep+\engwidth+\radius, 0);
\end{tikzpicture}
		\caption{}
		\label{fig:konwencja_znakow_1}
	\end{subfigure}
	\begin{subfigure}{0.45\textwidth}
		\centering
		\begin{tikzpicture}[x=0.5cm, y=0.5cm]
	\tikzmath{
		\tankwidth = 6;
		\tanksep = 5;
		\tankthickness = 1;
		\engwidth = 1;
		\engsep = 1;
		\radius = 2;
		\chansep = 1;
	}
	\draw[color=red, thick] (-\tankwidth, \tanksep) -- (\tankwidth, \tanksep);
	\draw[color=red, thick] (-\tankwidth, \tanksep) -- (-\tankwidth, \tanksep+\tankthickness);
	\draw[color=red, thick] (\tankwidth, \tanksep) -- (\tankwidth, \tanksep+\tankthickness);
	\fill[color=red, opacity=0.2] (-\tankwidth, \tanksep) rectangle (\tankwidth, \tanksep+\tankthickness);
	
	\draw[color=blue, thick] (-\tankwidth, -\tanksep) -- (\tankwidth, -\tanksep);
	\draw[color=blue, thick] (-\tankwidth, -\tanksep) -- (-\tankwidth, -\tanksep-\tankthickness);
	\draw[color=blue, thick] (\tankwidth, -\tanksep) -- (\tankwidth, -\tanksep-\tankthickness);
	\fill[color=blue, opacity=0.2] (-\tankwidth, -\tanksep) rectangle (\tankwidth, -\tanksep-\tankthickness);
	
	\draw[black] (-\engsep-\engwidth, -\tanksep) -- (-\engsep-\engwidth, \tanksep);
	\draw[black] (-\engsep+\engwidth, -\tanksep) -- (-\engsep+\engwidth, -\chansep-\radius);
	\draw[black] (-\engsep+\engwidth, \chansep+\radius) -- (-\engsep+\engwidth, \tanksep);
	\draw[black] (-\engsep+\engwidth, -\chansep-\radius) arc (180:90:\radius);
	\draw[black] (-\engsep+\engwidth, \chansep+\radius) arc (180:270:\radius);
	\node[anchor=east] at (-\engsep-\engwidth, 0.5*\tanksep) {\(\vab{Q_2}\)};
	\node[anchor=east] at (-\engsep-\engwidth, -0.5*\tanksep) {\(\vab{Q_1}\)};
	\node[anchor=west] at (-\engsep+\engwidth+\radius, 0) {\(\vab{W}\)};
	\draw[line width=0.5mm, -{Stealth[length=5mm, open]}] (-\engsep, 0.75*\tanksep) -- (-\engsep, 0.25*\tanksep);
	\draw[line width=0.5mm, -{Stealth[length=5mm, open]}] (-\engsep, -0.75*\tanksep) -- (-\engsep, -0.25*\tanksep);
	\draw[line width=0.5mm, -{Stealth[length=5mm, open]}] (-\engsep+\engwidth+\radius, 0) -- (-\engsep+\engwidth, 0);
\end{tikzpicture}
		\caption{}
		\label{fig:konwencja_znakow_2}
	\end{subfigure}
	\captionsetup{subrefformat=parens}
	\caption{Porównanie konwencji znaków: techniczna \subref{fig:konwencja_znakow_1} i statystyczna \subref{fig:konwencja_znakow_2}.}
\end{figure}

Z Pierwszej Zasady Termodynamiki wiadomo, że energia wewnętrzna \(U\) może być zmieniona jedynie poprzez wymianę ciepła \(Q\) lub wykonanie pracy nad układem \(W\) i ma formę:
\begin{equation*}
	\Delta U = Q + W,
\end{equation*}
co w formie różniczkowej ma postać \cite{b:huang_mechanika_statystyczna}:
\begin{equation}\label{eq:def_1zt}
	\d{U} = \idif{Q} + \idif{W},
\end{equation}
gdzie \(\dbar\) oznacza różniczkę niezupełną, czyli taką której całka zależy od drogi przemiany, wzdłuż której następuje całkowanie. W przeciwieństwie do niej \(\d{U}\) jest różniczką zupełną, co oznacza, że całka nie zależy od drogi, a \(U\) jest funkcją stanu. Całkując wyrażenie \eqref{eq:def_1zt} po jednym pełnym cyklu, pamiętając że \(\oint \d{U} = 0\), otrzymujemy:
\begin{equation}\label{eq:silnik_dU0}
	\oint \d{U} = 0 = \oint \idif{Q} + \oint \idif{W}, \quad\text{stąd}\quad
	-W = Q.
\end{equation}
Oznacza to, że suma ciepeł wymienionych między układem a otoczeniem jest równa pracy wykonanej przez układ \(-W\).

Silnik ma kontakt cieplny z dwoma zbiornikami o różnych temperaturach. Ciepło możemy rozdzielić na to dostarczone ze zbiornika ciepłego \(Q_h\) i na to oddane do zbiornika zimnego \(Q_c\):
\begin{equation}\label{eq:ciepla_w_silniku}
	Q = Q_h + Q_c.
\end{equation}
Zgodnie z równaniem \eqref{eq:silnik_dU0} praca wykonana jest równa:
\begin{equation}\label{eq:praca_w_silniku}
	-W = Q_h + Q_c.
\end{equation}
W cyklu prawobieżnym, w silnikach wykorzystujących gaz jako czynnik roboczy, praca wykonana \(W = - \oint p \d{V}\) jest kosztem ciepła \(Q_h\) pobranego ze zbiornika o temperaturze \(T_h\).

Parametrem opisujących efektywność przemiany dostarczonego do silnika ciepła na uzyskaną pracę jest sprawność silnika cieplnego:
\begin{equation}\label{eq:def_sprawnosc_silnika_W}
\eta = \frac{-W}{Q_h}.
\end{equation}
Biorąc pod uwagę relację \eqref{eq:praca_w_silniku} sprawność możemy zapisać poprzez ciepła \(Q_h\) i \(Q_c\):
\begin{equation}\label{eq:def_sprawnosc_silnika}
	 \eta = \frac{Q_h + Q_c}{Q_h} = 1 + \frac{Q_c}{Q_h}.
\end{equation}
Jest ona wielkością ważną zarówno w termodynamice teoretycznej jak i technicznej. W silnikach chcemy uzyskać jak największy zysk w postaci pracy jak najmniejszym kosztem w postaci dostarczonego ciepła \cite{b:pudlik_termodynamika}.

W analizie wyidealizowanych silników cieplnych rozważa się je jako serię elementarnych przemian: adiabatycznej (bez wymiany ciepła), izotermicznej (o stałej temperaturze czynnika roboczego), izobarycznej (o stałym ciśnieniu) oraz izochorycznej (o stałej objętości). Łącząc te przemiany w cykle zamknięte możemy otrzymać różnorodne cykle silników cieplnych o różnej liczbie przemian od 3, poprzez najpopularniejsze o 4 przemianach, i o większej ich liczbie. Przeanalizowane zostaną tu cykle: Carnota, Otta, Stirlinga.

\subsection{Cykl Carnota}
Silnik Carnota jest serią czterech przemian, w którym czynnikiem roboczym jest gaz doskonały spełniający równanie stanu:
\begin{equation}\label{eq:def_gaz_doskonaly}
	pV = nRT,
\end{equation}
gdzie \(p\) to ciśnienie gazu doskonałego, \(V\) -- objętość zajmowana przez ten gaz, \(n\) -- liczba moli tego gazu, \(R \approx \qty{8.314}{\joule / \kelvin / \mol}\) -- uniwersalna stała gazowa, \(T\) -- temperatura gazu.

\begin{figure}
	\centering
	\begin{tikzpicture}
	\begin{axis}[width=0.7\textwidth,
		xlabel=$V$,
		ylabel=$p$,
		axis x line=bottom,
		axis y line=middle,
		xmin=0,
		xmax=1.2,
		ymin=0,
		ymax=1.2,
		ticks=none,
		every axis x label/.style={at={(current axis.right of origin)},anchor=north east},
		every axis y label/.style={at={(current axis.above origin)},anchor=north east},
		declare function={
			r = 2;
			tau = 2;
			f = 3;
			kappa = 1 + 2/f;
			V1 = 1;
			V2 = 1/r;
			V3 = 1/r * tau^(-1/(kappa-1));
			V4 = tau^(-1/(kappa-1));
			p1 = 1/r * tau^(-kappa/(kappa-1));
			p2 = tau^(-kappa/(kappa-1));
			p3 = 1;
			p4 = 1/r;
		}]
		\addplot[mark=*] coordinates {(V1,p1)} node[anchor=south west]{$1$};
		\addplot[mark=*] coordinates {(V2,p2)} node[anchor=north east]{$2$};
		\addplot[mark=*] coordinates {(V3,p3)} node[anchor=south east]{$3$};
		\addplot[mark=*] coordinates {(V4,p4)} node[anchor=south west]{$4$};
		\addplot[domain=V2:V1] {1/r * tau^(-kappa/(kappa-1)) / x}             node [pos=0.3, below, sloped] {$T_c$};
		\addplot[domain=V3:V2] {1/r^kappa * tau^(-kappa/(kappa-1)) / x^kappa} node [pos=0.5, below, sloped] {$Q=0$};
		\addplot[domain=V3:V4] {1/r * tau^(-1/(kappa-1)) / x}                 node [pos=0.5, above, sloped] {$T_h$};
		\addplot[domain=V4:V1] {1/r * tau^(-kappa/(kappa-1)) / x^kappa}       node [pos=0.5, above, sloped] {$Q=0$};
		\draw[line width=0.3mm, -{Stealth[length=5mm, open]}] (0.4,0.7) node[anchor=south west]{$Q_{3-4}$} -- (0.3,0.6);
		\draw[line width=0.3mm, -{Stealth[length=5mm, open]}] (0.75,0.1) -- (0.75,0.025) node[anchor=west]{$Q_{1-2}$};
	\end{axis}
\end{tikzpicture}

	\caption{Wykres cyklu Carnota w obiegu prawobieżnym.}
	\label{fig:pv_silnik_carnot}
\end{figure}

Przemiany w cyklu Carnota zaprezentowane są na rysunku \ref{fig:pv_silnik_carnot}. Są to po kolei: sprężanie izotermiczne w temperaturze \(T_c\); sprężanie adiabatyczne, której skutkiem jest podniesienie temperatury do \(T_h\); rozprężanie izotermiczne w tej samej temperaturze; oraz rozprężanie adiabatyczne do temperatury początkowej. Każda z przemian silnika Carnota jest odwracalna, a więc i cały cykl jest odwracalny. Przemiany izotermiczne występują bez różnicy temperatur między czynnikiem roboczym a zbiornikiem cieplnym. Energia wewnętrzna gazu doskonałego zależy jedynie od temperatury \cite{b:kondepudi_prigogine_modern_thermodynamics}. Oznacza to, że podczas przemian izotermicznych \(\Delta U = 0\), a \(- W = Q\).

Silnik taki można uruchomić w drugą stronę. Spowoduje to zmianę znaków \(W\), \(Q_{3-4}\) oraz \(Q_{1-2}\) i w efekcie kosztem włożonej pracy pobrane zostanie ciepło \(Q_{1-2}\) z zimniejszego zbiornika o temperaturze \(T_c\) i oddanie ciepła \(Q_{3-4}\) do cieplejszego o temperaturze \(T_h\). Aby przemiana izotermiczna była odwracalna musi być nieskończenie powolna. Cykl ten jest więc nieosiągalny w rzeczywistości, ale ma duże znaczenie w teorii termodynamiki, ponieważ stanowi górną granicę dla sprawności silników cieplnych.

W silniku Carnota ciepła oddane do zbiornika o temperaturze \(T_c\) i dostarczone ze zbiornika o temperaturze \(T_h\) są równe odpowiednio:
\begin{gather*}
	Q_{1-2} = - W_{1-2} = \int_{V_1}^{V_2} p \d{V} = nRT_c \ln\pab{\frac{V_2}{V_1}}, \\
	Q_{3-4} = - W_{3-4} = \int_{V_3}^{V_4} p \d{V} = nRT_h \ln\pab{\frac{V_4}{V_3}},
\end{gather*}
gdzie \(V_1\), \(V_2\), \(V_3\) oraz \(V_4\) to objętości gazu w silniku w punktach oznaczonych na rysunku \ref{fig:pv_silnik_carnot}. Objętości te są powiązane ze sobą warunkami wynikającymi z przemian adiabatycznych \cite{b:pudlik_termodynamika}:
\begin{align*}
	p_{2}V_{2}^{\varkappa} &= p_{3}V_{3}^{\varkappa}, \\
	p_{1}V_{1}^{\varkappa} &= p_{4}V_{4}^{\varkappa},
\end{align*}
gdzie parametr \(\varkappa\) to wykładnik adiabaty zależny od użytego gazu. Parametr ten wyraża się przez ciepło molowe przy stałym ciśnieniu \(c_p\) i ciepło molowe przy stałej objętości \(c_V\) wzorem \(\varkappa = \frac{c_p}{c_V}\). Wykorzystując wzór \eqref{eq:def_gaz_doskonaly} możemy znaleźć:
\begin{equation*}
	\frac{V_1}{V_2} = \frac{V_4}{V_3}.
\end{equation*}
Sprawność silnika jest równa:
\begin{equation}\label{eq:sprawnosc_silnika_Carnota}
	\eta = 1 + \frac{Q_{1-2}}{Q_{3-4}} =
	1 + \frac{nRT_c \ln\pab{\frac{V_2}{V_1}}}{nRT_h \ln\pab{\frac{V_4}{V_3}}} =
	1 - \frac{T_c}{T_h}.
\end{equation}
Wzór ten pokazuje, że sprawność silnika Carnota jest funkcją jedynie stosunku temperatur zbiorników.

\subsection{Cykl Otto}
Cykl Otto składa się ze sprężania adiabatycznego (\(1 \rightarrow 2\)), podgrzania izochorycznej (\(2 \rightarrow 3\)), rozprężania adiabatycznego (\(3 \rightarrow 4\)) oraz ochłodzenia izochorycznego (\(4 \rightarrow 1\)). Przedstawiono go na rysunku \ref{fig:pv_silnik_otto}. Silnik ten pracuje pomiędzy dwoma zbiornikami o różnych temperaturach. Ciepło dostarczane jest ze zbiornika o temperaturze \(T_h\), a oddawane do zbiornika o temperaturze \(T_c\). Podczas pobierania ciepła ze zbiornika o temperaturze \(T_h\) czynnik roboczy znajduje się w temperaturze niższej niż \(T_h\). Natomiast podczas oddawania ciepła do drugiego zbiornika o temperaturze \(T_c\) znajduje się on w temperaturze od niego wyższej. W trakcie tych przemian temperatura zbiorników jest stała i przepływy ciepła nie mają na nie wpływu. Przemiany te zachodzą więc przy spadku temperatury na ściankach silnika, co czyni go nieodwracalnym. Zgodnie z \eqref{eq:praca_w_silniku} aby określić pracę wykonaną w cyklu należy znaleźć ciepło dostarczone oraz odprowadzone, co ma miejsce jedynie podczas przemian (\(2 \rightarrow 3\)) oraz (\(4 \rightarrow 1\)). W przemianie izochorycznej według wzoru \eqref{eq:def_1zt} ciepło to jest równe zmianie energii wewnętrznej i jest równe \cite{b:pudlik_termodynamika}:
\begin{equation}
	\idif{Q} = \pab{\pdv{U}{T}}_{V} \d{T} = n c_{V} \d{T},
\end{equation}
gdzie \(c_{V}\) to ciepło molowe dla stałej objętości, które dla gazu doskonałego jest parametrem niezależnym od temperatury. Całkując to wyrażenie podczas przemian otrzymujemy:
\begin{equation}\label{eq:otto_Q}
\begin{split}
	Q_{2-3} = n c_{V} \pab{T_3 - T_2}, \\
	Q_{4-1} = n c_{V} \pab{T_1 - T_4},
\end{split}
\end{equation}
gdzie temperatury \(T_1\), \(T_2\), \(T_3\) oraz \(T_4\) to temperatury w punktach oznaczonych na wykresie \ref{fig:pv_silnik_otto}. Sprawność tego cyklu wyraża się wzorem:
\begin{equation}
	\eta = 1 + \frac{Q_{2-3}}{Q_{4-1}} = 1 - \frac{T_3 - T_2}{T_4 - T_1}.
\end{equation}
Wzór ten jest zależny od temperatur, których nie kontrolujemy bezpośrednio. Należy go wyrazić poprzez inne wielkości.

\begin{figure}
	\centering
	\input{figures/pv_silnik_otto.tex}
	\caption{Wykres cyklu Otto w obiegu prawobieżnym.}
	\label{fig:pv_silnik_otto}
\end{figure}

Parametry gazu w przemianie adiabatycznej spełniają wyrażenia pomocnicze \cite{b:pudlik_termodynamika}:
\begin{equation}
\begin{split}
	p_{1}V_{1}^{\varkappa} &= p_{2}V_{2}^{\varkappa}, \\
	p_{4}V_{1}^{\varkappa} &= p_{3}V_{2}^{\varkappa}.
\end{split}
\end{equation}
Obie z powyższych równości możemy przekształcić, aby ciśnienia zależały od stosunku objętości:
\begin{equation*}
	\frac{p_2}{p_1} = \frac{p_3}{p_4} = \pab{\frac{V_1}{V_2}}^{\varkappa} = r^{\varkappa}.
\end{equation*}
Wykorzystując równanie stanu gazu doskonałego \eqref{eq:def_gaz_doskonaly} możemy znaleźć analogiczną formę zależną of temperatury i objętości:
\begin{equation}
\begin{split}
	T_{1}V_{1}^{\varkappa-1} &= T_{2}V_{2}^{\varkappa-1}, \\
	T_{4}V_{1}^{\varkappa-1} &= T_{3}V_{2}^{\varkappa-1}.
\end{split}
\end{equation}
Dzieląc stronami otrzymujemy stosunki temperatur. Podobnie jak wyżej możemy wyrazić stosunek temperatur też za pomocą współczynnika sprężenia:
\begin{gather}
	\frac{T_3}{T_2} = \frac{T_4}{T_1}, \\
	\frac{T_2}{T_1} = \frac{T_3}{T_4} = \pab{\frac{V_1}{V_2}}^{\varkappa-1} = r^{\varkappa-1} \label{eq:otto_stosunek_T}.
\end{gather}
Posługując się tymi wielkościami możemy wyrazić temperatury \(T_2\) i \(T_3\) poprzez \(T_1\), \(T_4\), \(r\) i \(\varkappa\):
\begin{gather}
	T_2 = T_1 r^{\varkappa-1}, \\
	T_3 = T_4 r^{\varkappa-1}.
\end{gather}
Podstawiając je do wzorów na pracę oraz sprawność otrzymujemy:
\begin{gather}
	- W = n c_V \pab{r^{\varkappa-1} - 1} \pab{T_4 - T_1}, \label{eq:parametry_otto_W} \\
	\eta = 1 - \inv{r^{\varkappa-1}}. \label{eq:parametry_otto_eta}
\end{gather}
Możemy zauważyć, że sprawność nie zależy od temperatury a jedynie od stopnia sprężenia \(r\), który jest parametrem mechanicznym danego silnika. Ponieważ \(\varkappa > 1\) sprawność rośnie wraz ze wzrostem stopnia sprężenia \(r\).

\subsection{Cykl Stirlinga}
Podobnie jak cykl Carnota i cykl Otto, cykl Stirlinga składa się z czterech przemian. Są to po kolei: sprężanie izotermiczne w temperaturze \(T_c\) (\(1 \rightarrow 2\)), podgrzewanie izochoryczne (\(2 \rightarrow 3\)), rozprężenie izotermiczne w temperaturze \(T_h\) (\(3 \rightarrow 4\)) oraz ochłodzenie izochoryczne (\(4 \rightarrow 1\)). Cykl ten wyróżnia się dodatkowym elementem, którym jest wewnętrzny wymiennik ciepła. Jego rolą jest częściowy odzysk ciepła odprowadzanego podczas jednej z przemian izochorycznych i następnie doprowadzenie go podczas drugiej z nich. Otrzymuje się dzięki temu wyższą sprawność całego cyklu. Odpowiednio ciepła w każdej z przemian są równe:
\begin{gather}
	Q_{1-2} = - W_{1-2} = \int_{V_1}^{V_2} p \d{V} = nRT_c \ln\pab{\frac{V_2}{V_1}} \\
	Q_{2-3} = n c_V \pab{T_h - T_c} \\
	Q_{3-4} = - W_{3-4} = \int_{V_2}^{V_1} p \d{V} = nRT_h \ln\pab{\frac{V_1}{V_2}} \\
	Q_{4-1} = n c_V \pab{T_c - T_h}.
\end{gather}

\begin{figure}
	\centering
	\input{figures/pv_silnik_stirling.tex}
	\caption{Wykres cyklu Stirlinga w obiegu prawobieżnym.}
	\label{fig:pv_silnik_stirling}
\end{figure}

Ciepło \(Q_{2-3}\) dostarczone podczas przemiany izochorycznej jest dokładnie równe ciepłu \(Q_{4-1}\) oddanemu w odpowiadającej przemianie, ale z przeciwnym znakiem. Ciepło to nigdy nie opuszcza całego układu i jest jedynie wymieniane między gazem roboczym a regeneratorem. Praca tego cyklu, zgodnie ze wzorem \eqref{eq:praca_w_silniku}, ma postać:
\begin{equation}
	- W = n R \pab{T_h - T_c} \ln\pab{\frac{V_1}{V_2}},
\end{equation}
a sprawność zgodnie z \eqref{eq:def_sprawnosc_silnika_W} przy \(Q_h = Q_{3-4}\) jest równa:
\begin{equation}
	\eta = 1 - \frac{T_c}{T_h}.
\end{equation}
Sprawność tego cyklu jest identyczna ze sprawnością silnika Carnota, więc (co wynika z twierdzenia Carnota) jest on również w pełni odwracalny.